\chapter{Environment}
\section{Get Rolling}
\textbf{Look at the picture carefully.}
\vspace{1.4in}
\begin{enumerate}
    \item What is making the air dirty?
    \item List some other activities that can make air dirty.
    \item Discuss how you can contribute towards keeping the air clean.
\end{enumerate}
\section{What to Know}
Environment and its importance, causes, effects and ways of reducing different types of pollution, greenhouse effect and global warming, and waste management
\section{\textbf{Learning Objectives}}
\begin{itemize}
    \item to define environment and name its components
    \item to list the pollutants air, water and land
    \item to describe the causes and effects of different types of pollution
    \item to explain the ways for reducing different types of pollution
    \item to apply the 3RS of waste management in real life
\end{itemize}
\section{I Explore}
All the things around us that we see and feel make up our environment. The air we breathe, the water we drink, the plants, trees and animals around us and even weather-all constitute the environment. In other words, the natural surroundings in which an organism lives is called its \textbf{environment.}
\section{Environment and its Components}
The environment is made up of both living and non-living components. Plants, animals and human beings are the living components of the environment. Air, water, soil, light and temperature are the non-living components of the environment.
\subsection{Importance of the Environment to Us and Other living Beings}
The environment is essential for life on Earth. It provides us with everything we need to live-clean air to breathe, fresh water to drink and food to eat. It also regulates the temperature of our surroundings and brings about a change in the climate, which is crucial for the survival of all living beings.
\vspace{0.1in}
Taking care of the environment is crucial for our well-being of our planet.
\subsection{Impact of Human Activities on the Environment}
The number of people living on Earth is referred to as its \textbf{human population.} Human population is increasing day by day. This has led to an increase in the demand for food, water, housing and clothing. Hence, forests are cut down to get more and more clear land for building houses and factories. Some human activities have deep impacts on the environment. For example, the population rise has led to an increased vehicle use, which has resulted in air pollution. Improper waste disposal has caused soil pollution, while waste disposal in water bodies has resulted in water pollution. The machines used in factories create a lot of noise that leads to noise pollution. Let us learn more about pollution and its types.
\section{Pollution and its Types}
\textbf{Pollution} refers to the introduction of harmful substances or contaminants into the environment, causing negative effects on the natural world and often on human health. These harmful substances, known as \textbf{pollutants,} can come from various sources, including human activities and natural processes.
\vspace{0.1in}
\par
There are several types of pollution, each with is own specific causes and consequences. Some of the types are discussed here.
\subsection{Air Pollution}
This occurs when harmful substances are released into the air, leading to the deterioration of air quality. Common air pollutants include carbon monoxide, sulphur dioxide, nitrogen oxides and particulate matter. Air pollution can cause smog and breathing difficulties and can also damage plant and animal life.
\subsection{\textbf{Causes of air pollution}}
Some major causes of air pollution include vehicular and industrial exhaust. Additionally, burning of fossil fuels, deforestation and dust and debris generated from construction and demolition activities contribute to air pollution.
\subsection{\textbf{Effects of air pollution}}
Air pollution has a wide range of negative effects on both the environment and human health. For example, in humans, pollution can cause respiratory disorders like asthma and cough. Over time, it can cause serious lung diseases too. Air pollution caused by the release of some gases like carbon dioxide in the air has led to global warming and climate change, leading to shifts in weather patterns, rising sea levels, etc.
\vspace{0.1in}
\par
I wonder what would happen if there were no greenhouse gases at all in the atmosphere!
\vspace{0.1in}
\par
But what is global warming? Earth has an atmosphere that is made of gases such as nitrogen, oxygen and carbon dioxide. Certain gases let the Sun's heat get through the atmosphere towards the earth, bout do not let it escape back. These gases are called \textbf{greenhouse gases.} Why are they called greenhouse gases?
\vspace{0.1in}
\par
A greenhouse is where plants are grown in an enclosed space with walls made of glass. Glass allows sunrays to enter the glasshouse but it does not allow the heat to escape.
\vspace{0.4in}
\par
Therefore, the heat of the Sun warms the inside of the glass use and also makes the air moist, which is suitable for plant growth.
\vspace{0.1in}
\par
A similar effect is caused by some of the gases in the atmosphere. Therefore, these gases are called \textbf{greenhouse gases.}
\vspace{0.1in}
\par
One of the greenhouse gases is carbon dioxide. Now, if the level of carbon dioxide around the earth increases, the earth gets warmer. This is called \textbf{global warming.} Global warming has led to many effects on weather, sea level and habitats of animals.
\vspace{0.1in}
\par
The ice on mountains and at the North and the South Poles can melt, which can make sea levels go up. Some animals might not have places to live because their habitats are getting too warm. For example, penguins, sea turtles, etc. are facing the danger of lossing habitats.
\vspace{2.4in}
\subsection{Water Pollution}
\textbf{Water pollution} when pollutants are released into water bodies like rivers, lakes and oceans.
\subsection{\textbf{Causes of water pollution}}
Sometimes big factories and machines release harmful chemicals into water, making it unfit for consumption. Other causes include dirty water from homes, oil spills and trash such as plastic bottles, bags and other garbage that people throw into water bodies.
\subsection{\textbf{Effects of water pollution}}
Water pollution causes a number of diseases like typhoid, dysentery and cholera. These diseases take millions of lives every year. Fish and other marine animals are also adversely affected by water pollution.
\subsection{Soil Pollution}
\textbf{Soil pollution} happens when harmful things contaminate the soil where plants grow.
\subsection{\textbf{Causes of soil pollution}}
Soil pollution involves the contamination of soil with pollutants such as heavy metals, farm chemicals like pesticides and industrial chemicals. Heavy metals enter the soil through mining activities and e-waste. Farmers often use chemicals to promote plant growth, and when these chemicals are used excessively, they enter the soil with water, leading to pollution. Another significant cause of soil pollution is improper waste disposal. If people do not dispose of their waste properly but just leave it on the ground, it can make the soil dirty and smelly.
\vspace{0.2in}
\par
I wonder how improper waste disposal contributes to pollution!
\subsection{\textbf{Effects of soil pollution}}
Polluted soil can make it difficult for plants to grow as they might not receive the necessary nutrients. This can lead to a reduction in agricultural productivity and animals dependent on soil. Moreover, the chemicals that enter the soil may find their way into the food we eat, causing various health hazards.
\subsection{Noise Pollution}
\textbf{Noise pollution} is excessive, unwanted or disruptive noise that can have negative effects on human health.
\subsection{\textbf{Causes of noise pollution}}
Noise pollution is caused by loud sounds from vehicles, machines, construction sites and music that disturb the environment an people's peace. Loud noise from activities like parties, airplanes and factories can also contribute to noise pollution.
\subsection{\textbf{Effect of noise pollution}}
Continuous exposure to noise pollution can lead to serious health issues. It can damage our ears and, over time, lead to hearing loss. Additionally, it can cause stress, anxiety and even affect attention and mental well-being. Noise pollution is also linked to heart problems and many other health issues.
\vspace{1.6in}
\section{Saving the Environment}
We can work together to reduce pollution in different ways. Here are some simple steps to reduce each type of pollution.
\subsection{Air Pollution}
\begin{itemize}
    \item Use public transport, walk or bike instead of driving alone. This would reduce the number of vehicles on the road.
    \item Vehicles must be checked regularly for pollution.
    \item Plant trees around your home to improve air quality. Plants help reduce the carbon dioxide content in the atmosphere.
    \item Avoid burning trash. Burning produces a lot of smoke, adding a lot of harmful gases to the environment.
\end{itemize}
\subsection{Water Pollution}
\begin{itemize}
    \item Don't throw garbage or plastics into rivers, lakes, oceans or other waterbodies.
    \item Use fewer plastic bags. Use recycled bags instead.
    \item Properly dispose of chemicals, paints and medicines to keep them out of water.
\end{itemize}
\subsection{Soil Pollution}
\begin{itemize}
    \item Reduce, reuse and recycle to create less waste.
    \item Use natural substances for gardening and farming.
    \item Do not throw e-waste in the open.
    \item Avoid using harmful chemicals in your garden or yard.
\end{itemize}
\subsection{Noise Pollution}
\begin{itemize}
    \item Keep the volume of music and TV to a comfortable level.
    \item Be mindful of the noise you make when playing outdoors.
    \item Use headphones when listening to music or watching videos in public places.
    \item Do not honk until absolutely necessary.
    \item Follow rules and guidelines for limiting noise in your community.
\end{itemize}
Another important way for saving the environment is proper waste management. Let us study more about it.
\subsection{Waste Management}
\textbf{Waste management} is about managing waste and making sure it does not harm the environment. There are three important things we can do to manage waste properly. We called them the \textbf{3RS.}
\vspace{0.1in}
\par
\textbf{Reduce:} This refers to producing less waste in the first place. We can do this by using stuff wisely; for example, by using both sides of a paper before throwing it away.
\vspace{0.1in}
\par
\textbf{Reuse:} Instead of throwing things away after using  the things once, we can find new ways to use them again. For example, using an old jar to store things or using old clothes as cleaning rags.
\vspace{0.1in}
\par
\textbf{Recycle:} In recycling, we send certain things, like paper, plastic and glass, to special places where they can be turned into new things. It is like giving them a second chance to be useful.
\vspace{0.1in}
\par
By following the 3RS (reduce, reuse and recycle), we can help keep our planet clean and protect it for future generations.
\vspace{3.2in}
\section{Now I Know}
\begin{itemize}
    \item The natural surroundings in which organism lives is called its \textbf{environment.} It consists of both living and non-living components.
    \item \textbf{Pollution} refers to the introduction of harmful substances or \textbf{contaminants} into the environment, causing negative effects on the natural world and often on human health.
    \item \textbf{Air pollution} occurs when harmful substances are released into the air, leading to the deterioration of air quality.
    \item Air pollution caused by the release of some gases like carbon dioxide has led to global warming and climate change, leading in turn to shifts in weather patterns, rising sea levels, etc.
    \item \textbf{Water pollution} happens when pollutants are released into waterbodies like rivers, lakes and oceans.
    \item \textbf{Soil pollution} involves the contamination of soil with pollutants like heavy metals, farm chemicals like pesticides and industrial chemicals.
    \item \textbf{Noise pollution} is caused by loud sounds from vehicles, machines, construction activities and music that disturb the environment and people's peace.
    \item Important ways for saving the environment include proper waste management and controlling pollution through different ways.
    \item The \textbf{3RS} for waste management are reduce, reuse and recycle.
\end{itemize}